\scanpage{91}

\[
(28)\quad
\left\{\begin{aligned}
\widetilde{\Pi}_{0,3}=\Pi_{0,3}\cdot\{1,\tau\}&\longrightarrow\widetilde{\Pi}'_{0,3}\\
\tau_0,\tau_1,\tau_\infty&\longmapsto
\sigma'_0,\sigma'_1,\sigma'_\infty
\end{aligned}\right.
\]

This being so, we define an \uline{isomorphism} \deleted{\illegible{1}}
\[
(29)\quad
\begin{aligned}
\widetilde{\Pi}'_{0,3}&\xrightarrow{\sim}\mathfrak S_{0,3}\\
(\sigma'_0,\sigma'_1,\sigma'_\infty)
&\longrightarrow
\left(\deleted{\unsure{\sigma^{\prime}_1}}
\underset{\substack{\inserted{\text{over }\dot\tau'_\infty}\\\Vert\\\inserted{\dot\sigma_\infty\tau}}}{\tau'_\infty},
\underset{\substack{\inserted{\text{over }\dot\tau'_0}\\\Vert\\\inserted{\dot\sigma_0\tau}}}{\tau'_0},
\underset{\inserted{\text{over }\tau}}{\tau_\infty}\right)
\end{aligned}
\]
also denoted $r_0,r_1,r_\infty$ (fundamental \uline{reflections})

\margin{This isomorphism is associated with a canonical homomorphism
$\varphi:\widetilde{\Pi}_{0,3}=\Pi_{0,3}\cdot\{1,\tau\}\xrightarrow{\text{epimorphism}}\mathfrak S_{0,3}$
inducing a canonical homomorphism
$\varphi:\Pi_{0,3}\xrightarrow{\text{epimorphism}}\mathfrak S_{0,3}^{+}$
allowing us to identify $\mathfrak S_{0,3}$ and $\mathfrak S_{0,3}^{+}$
with the quotients of the left-hand~\unsure{sides|sides} by the relations
$\ell_1^2=\ell_\infty^3=1$ \ldots}

which induces an isomorphism
\[
(30)\quad
\left\{\begin{aligned}
\Pi'_{0,3}&\xrightarrow{\sim}\deleted{\unsure{\Pi}}\mathfrak S_{0,3}^{+}\\
\rho_s=\ell'_0&\longmapsto\varepsilon_0=\sigma_\infty\rho\\
\sigma=\ell'_1&\longmapsto\sigma_\infty\\
\rho_f=\ell'_\infty&\longmapsto\rho^{-1}
\end{aligned}\right.
\]

\margin{\uline{N.B.}
$\left\{\begin{aligned}
\tau_\infty\tau'_0&=\varepsilon_0=\ell'_0\\
\tau'_\infty\tau_\infty&=\sigma_\infty=\ell'_1\\
\tau'_0\tau'_\infty&=\rho^{-1}=\ell'_\infty
\end{aligned}\right.$
i.e.
$\begin{aligned}
r_\infty r_1&=\ell'_0\;(=\varepsilon_0)\\
r_0r_\infty&=\ell'_1\;(=\sigma_\infty)\\
r_1r_0&=\ell'_\infty\;(=\rho^{-1})
\end{aligned}$}

We thus find that $\mathfrak S_{0,3}$ has \deleted{\unsure{gen}} generating relations
\[
(31)\quad
{\tau'_\infty}^{2}={\tau'_0}^{2}=\tau_\infty^2
=\bigl(\underbrace{\tau'_\infty\tau_\infty}_{\sigma_\infty}\bigr)^2
=\bigl(\underbrace{\tau'_0\tau'_\infty}_{\rho^{-1}}\bigr)^3=1,
\]
and $\mathfrak S_{0,3}^{+}$ has generating relations
\[
(32)\quad \ell'_\infty \ell'_1\ell'_0={\ell'_1}^{2}={\ell'_\infty}^{3}=1
\]
or also simply, eliminating $\ell'_0$ and setting
$\ell'_1=\sigma_\infty$ and ${\ell'_\infty}^{-1}=\rho$
\[
(33)\quad\boxed{\sigma_\infty^2=\rho^3=1}
\]

For the record, we are going to express the notable elements of $\mathfrak S_{0,3}^{+}$ in terms of $\sigma_\infty,\rho$:

\scanpage{92}
\mspage{335}

\[
(34)\quad\left\{\begin{aligned}
\varepsilon_0&=\sigma_\infty\rho\ \deleted{\illegible{1}}&
\ell_0&=\sigma_\infty\rho\sigma_\infty\rho=(\sigma_\infty\rho)^2\\
\varepsilon_1&=\rho\sigma_\infty&
\ell_1&=\rho\sigma_\infty\rho\sigma_\infty=(\rho\sigma_\infty)^2\\
\varepsilon_\infty&=\rho^{-1}\sigma_\infty\rho^{-1}&
\ell_\infty&=(\rho^{-1}\sigma_\infty\rho^{-1})^2
=\rho^{-1}\sigma_\infty\rho^{-1}\sigma_\infty\rho^{-1}
\end{aligned}\right.,
\]
\[
(35)\quad\left\{\begin{aligned}
\sigma_0&=\rho^{\deleted{\unsure{2}}}\sigma_\infty\rho^{\deleted{\unsure{2}}-1}
\quad(=\deleted{\unsure{\rho}}\varepsilon_1\rho^{-1}=\rho\varepsilon_1^{-1})\\
\sigma_1&=\rho^{-1}\sigma_\infty\rho^{\deleted{\unsure{2}}}
\quad(=\rho^{-1}\varepsilon_0\deleted{\unsure{\rho}}=\varepsilon_\infty\rho^{-1})
\end{aligned}\right.
\]

and in $\mathfrak S_{0,3}^{-}$ we have, in terms of $\tau'_\infty,\tau'_0,\tau_\infty$,
the expressions for \deleted{$\unsure{\tau_0,\tau_1,\tau_\infty,\tau'_0,\tau'_1,\tau'_\infty}$}
\inserted{$\tau_0,\tau_1,(\tau_\infty,\tau'_0),\tau'_1,(\tau'_\infty)$}
\[
(36)\quad\left\{\begin{aligned}
\tau_0&=\deleted{\unsure{\tau_0\tau'_0\tau_\infty\tau'_\infty}}\rho\tau_\infty\rho^{-1}
=\tau'_\infty\tau'_0\tau_\infty\tau'_0\tau'_\infty
\quad(=r_0r_1r_\infty r_1r_0)\\
\tau_1&=\rho^{-1}\tau_\infty\rho
=\tau'_0\tau'_\infty\tau_\infty\tau'_\infty\tau'_0
\quad(=r_1r_0r_\infty r_0r_1)\\
\tau'_1&=\rho\tau'_0\rho^{-1}
=\tau'_\infty\deleted{\unsure{\tau'_0}}\tau'_0\tau'_\infty
\quad(=r_0r_1r_0)
\end{aligned}\right.
\]
and
\[
\left\{\begin{aligned}
\tau'_\infty&=r_0\\
\tau'_0&=r_1\\
\tau_\infty&=r_\infty.
\end{aligned}\right.
\]

\uline{IV}) Relation with
$\mathrm{Gl}(2,\mathbb Z)'=\mathrm{Gl}(2,\mathbb Z)/\mathord{\pm}1
\simeq\mathrm{GP}(1,\mathbb Z)$.

This group operates on the Poincaré half-plane
\[
D=\{z\in\mathbb C\mid\Im z>0\}
\]
by
\[
(37)\quad\left\{\begin{aligned}
u_g(z)&=\frac{az+b}{cz+d}&\text{if }ad-bc&=+1\\
u_g(z)&=\frac{a\bar z+b}{c\bar z+d}&\text{if }ad-bc&=-1
\end{aligned}\right.
\qquad
\begin{gathered}
g=\begin{pmatrix}a&b\\c&d\end{pmatrix}\\
a,b,c,d\in\mathbb Z\\
ad-bc\in\{\pm1\}
\end{gathered}
\]

We have in $D$ a hatched fundamental domain, which corresponds to the spherical
triangle $(-\bar j,0,\tfrac12)$ by the \unsure{homomorphism|homeomorphism}
of the natural \unsure{covering}, $i$ corresponding to $\tfrac12$,
$-\bar j$ to $-\bar j$, and the point at infinity to $0$.
\deleted{Thus \unsure{we find}}\inserted{The reflections}
with respect to the three sides of this triangle are given by
\[
(38)\quad
\deleted{\unsure{\tau_0}}\unsure{r_0|\tau_0}(z)=\frac1{\bar z},\qquad
\unsure{r_1|\tau_1}(z)=1-\bar z,\qquad
\unsure{r_\infty|\tau_\infty}(z)=-\bar z
\]
\figure{\begin{tikzpicture}[x=2.2cm,y=2.2cm,line cap=round,line join=round,font=\scriptsize]
\draw[line width=.4pt] (-1.55,0) .. controls (-.7,.04) and (.5,-.025) .. (1.85,.01);
\draw[line width=.4pt] (-1,0) arc[start angle=180,end angle=0,radius=1];
\draw[line width=.5pt] (0,-.48)--(0,2.85);
\draw[line width=.5pt] (.5,-.65)--(.5,2.78);
\begin{scope}
\clip (0,1) arc[start angle=90,end angle=60,radius=1]--(.5,2.78)--(0,2.85)--cycle;
\foreach \h in {.96,1.24,1.52,1.80,2.08,2.36,2.64}
  \draw[line width=.35pt] (0,\h)--(.5,{\h+.30});
\end{scope}
\draw[line width=.7pt] (0,1)--(.075,1)--(.075,1.075)--(0,1.075);
\fill (.5,.8660254) circle[radius=.025];
\node[inner sep=0pt] at (-1.04,-.14) {$-1$};
\node[inner sep=0pt] at (-.10,-.14) {$0$};
\node[inner sep=0pt] at (.43,-.18) {$\frac12$};
\node[inner sep=0pt] at (1.04,-.14) {$1$};
\node[inner sep=0pt] at (-.13,1.12) {$i$};
\node[inner sep=0pt] at (.64,.94) {$-\bar j$};
\node[inner sep=0pt,anchor=east,rotate=16] at (-.07,1.94) {$\dot r_\infty=\dot\tau_\infty$};
\node[inner sep=0pt,anchor=west] at (.56,2.25) {$\dot\tau'_0=\dot r_1$};
\node[inner sep=0pt,rotate=-63] at (.23,.63) {$\dot\tau'_\infty=\dot r_0$};
\end{tikzpicture}}

\deleted{\unsure{They determine}} and the pairwise products of these reflections

\scanpage{93}

are
\[
(39)\quad\left\{\begin{aligned}
\unsure{r_\infty r_1|\tau_\infty\tau_1}(z)&=\ell'_0(z)=z-1\;\deleted{\unsure{;}}\\\unsure{r_0r_\infty|\tau_0\tau_\infty}(z)&=\ell'_1(z)=-\frac1z\\
\unsure{r_1r_0|\tau_1\tau_0}(z)&=\ell'_\infty(z)=1-\frac1z=\frac{z-1}{z}
\end{aligned}\right.
\]

We thus find an isomorphism
\[
(40)\quad\left\{\begin{aligned}
\mathfrak S_{0,3}&\xrightarrow{\sim}\mathrm{Gl}(2,\mathbb Z)'\simeq
\mathrm{Gl}(2,\mathbb Z)/\mathord{\pm}1\simeq\mathrm{GP}(1,\mathbb Z)\\
r_0=\tau'_\infty&\longmapsto r_0=\begin{pmatrix}0&1\\1&0\end{pmatrix}\\
r_1=\tau'_0&\longmapsto r_1=\begin{pmatrix}-1&1\\0&1\end{pmatrix}\\
r_\infty=\tau_\infty&\longmapsto r_\infty=\begin{pmatrix}1&0\\0&-1\end{pmatrix}
\end{aligned}\right.
\]
\margin{it is better to change \unsure{some|of} signs \unsure{cf. Formulary|cf. formulary}}

and a presentation of $\mathrm{GP}(1,\mathbb Z)$ by
\[
(41)\quad r_0^2=r_1^2=r_\infty^2=(r_0r_\infty)^2=(r_1r_0)^3=1.
\]
This isomorphism induces an isomorphism
\[
(42)\quad\left\{\begin{aligned}
\mathfrak S_{0,3}^{+}&\xrightarrow{\sim}
\mathrm{Sl}(2,\mathbb Z)'=\mathrm{Sl}(2,\mathbb Z)/\mathord{\pm}1
\simeq\mathrm{GP}(1,\mathbb Z)^{+}\\
\ell'_1=\sigma_\infty&\longmapsto\begin{pmatrix}0&-1\\1&0\end{pmatrix}\\
\ell'_\infty=\rho^{-1}&\longmapsto\begin{pmatrix}1&-1\\1&0\end{pmatrix}
\end{aligned}\right.
\]
\[
(\text{hence}\quad\rho\longmapsto\begin{pmatrix}0&1\\-1&1\end{pmatrix}).
\]

We henceforth identify the elements of $\mathfrak S_{0,3}$ with their images in
$\unsure{\mathrm{Gl}|\mathrm{Sl}}(2,\mathbb Z)'\simeq\mathrm{GP}(1,\mathbb Z)$.
Let us note the following formulary \unsure{(where we use} that for a matrix
$u=\begin{pmatrix}a&b\\c&d\end{pmatrix}\in\mathrm{Sl}(2,\mathbb Z)$,
\deleted{\illegible{1}}
\deleted{\unsure{for} $\det u\in\mathbb Z^*=\{\pm1\}$ so $(\det u)^2=1$}, we have

\scanpage{94}
\mspage{336}

\[
(43)\quad u^{-1}=\begin{pmatrix}d&-b\\-c&a\end{pmatrix}\quad):
\]

\margin{N.B. \unsure{2.} \deleted{\unsure{If det}} $u\in\mathrm{Gl}(2,A)$,
$A$ an \unsure{arbitrary} ring, we have
$\det(u)\cdot u^{-1}=\begin{pmatrix}d&-b\\-c&a\end{pmatrix}$}

\deleted{\unsure{And}}.
\[
(44)\quad\left\{\begin{aligned}
\varepsilon_0=\sigma_\infty\rho&=\begin{pmatrix}1&-1\\0&1\end{pmatrix},&
\varepsilon_0^{-1}&=\begin{pmatrix}1&1\\0&1\end{pmatrix}\\
\varepsilon_1=\rho\sigma_\infty&=\begin{pmatrix}1&0\\1&1\end{pmatrix},&
\varepsilon_1^{-1}&=\begin{pmatrix}1&0\\-1&1\end{pmatrix}\\
\varepsilon_\infty=\rho^{-1}\sigma_\infty\rho^{-1}&=\begin{pmatrix}2&-1\\1&0\end{pmatrix},&
\varepsilon_\infty^{-1}&=\begin{pmatrix}0&1\\-1&2\end{pmatrix}\\
\ell_0=\varepsilon_0^2&=\begin{pmatrix}1&-2\\0&1\end{pmatrix}&
\sigma_0&=\begin{pmatrix}1&-1\\2&-1\end{pmatrix}\\
\ell_1=\varepsilon_1^2&=\begin{pmatrix}1&0\\2&1\end{pmatrix}&
\sigma_1&=\begin{pmatrix}1&-2\\1&-1\end{pmatrix}\\
\ell_\infty=\varepsilon_\infty^2&=\begin{pmatrix}3&-2\\2&-1\end{pmatrix}&
\sigma_\infty&=\begin{pmatrix}0&-1\\1&0\end{pmatrix}
\end{aligned}\right.
\]
\margin{Cf. ``formulary'' for a more \unsure{complete} list of relations and expressions.}
\[
(45)\quad\rho=\begin{pmatrix}0&1\\-1&1\end{pmatrix},\qquad
\unsure{\rho^{-1}|\rho^{-1\prime}}=\begin{pmatrix}1&-1\\1&0\end{pmatrix},
\]
\[
(46)\quad \ell_\infty \ell_1\ell_0=\begin{pmatrix}-1&0\\0&-1\end{pmatrix}
\]
the nontrivial element of the center of $\mathrm{Sl}(2,\mathbb Z)$ which
(fortunately!) gives $1$ in $\mathrm{Sl}(2,\mathbb Z)'$

The group $\Pi_{0,3}$ is recovered as the kernel of the homomorphism
\diagram{\begin{tikzcd}[ampersand replacement=\&]
\mathrm{Sl}(2,\mathbb Z)' \arrow[r] \& \mathrm{Sl}(2,\mathbb Z/2\mathbb Z) \arrow[d,phantom,"\simeq"{description,rotate=90}]\\
\& \mathfrak S_3
\end{tikzcd}}
the exact sequence
\[
(47)\quad
\diagram{\begin{tikzcd}[ampersand replacement=\&]
1 \arrow[r] \& \Pi_{0,3}\arrow[r,dash] \& \mathfrak S_{0,3}^{+}\arrow[r,dash] \& \mathcal{T}_{0,3}^{+\unsure{\prime}}\arrow[r,dash]\arrow[d,phantom,"\simeq"{description,rotate=90}] \& 1\\
\&\&\&\mathfrak S_3\&
\end{tikzcd}}
\]
is recovered as
\[
(48)\quad
\diagram{\begin{tikzcd}[ampersand replacement=\&]
1\arrow[r]\&\Pi(2)'\arrow[r]\arrow[d,phantom,"\simeq"{description,rotate=90}] \& \mathrm{Sl}(2,\mathbb Z)'\arrow[r] \& \mathrm{Sl}(2,\mathbb Z/2\mathbb Z)\arrow[r]\arrow[d,phantom,"\simeq"{description,rotate=90}] \& 1\\
\&\Pi_{0,3}^{\unsure{\beta}}\&\&\mathfrak S_3\&
\end{tikzcd}}
\]

\scanpage{95}

It makes \unsure{our hands|our hand} itch to \deleted{\unsure{fool around}}\inserted{work}
in $\mathrm{Gl}(2,\mathbb Z)$ itself, a central extension of
$\mathfrak S_{0,3}$ --- in fact, this extension is interpreted as
$\mathcal{T}_{1,1}$ --- but we shall come back to it, later\ldots

\uline{V}) Let us sum up the avatars of the same group, which we have
\unsure{encountered|encountered} so far

1) $\mathfrak S_{0,3}\overset{\text{dfn}}{=}\operatorname{Aut}_{\mathrm{lac}}(\Pi_{0,3})$
with \unsure{subgroup} $\mathfrak S_{0,3}^{+}$ (preserving orientation)

2) $\mathcal{T}'_{0,4}\simeq\operatorname{Autext}_{\mathrm{lac}}(\Pi_{0,4})'$
outer automorphisms of $\Pi_{0,4}$ respecting the loop structure and fixing the loop class
$\unsure{L_4|L_\lambda}$ (and permuting $L_0,L_1,\unsure{L_\infty|L_3}$)

\deleted{with $\unsure{\mathcal{T}_{0,4}^{\prime+}}$}
\deleted{$\unsure{A'_{0,4}/A_{0,4}^{\prime0}}\simeq\unsure{\operatorname{Auttop}}(\mathbb P^1_{\mathbb C}\setminus\{0,\unsure{1}\})$}
with $\mathcal{T}_{0,4}^{\prime+}$ (preserving orientation)

3) $A'_{0,4}/A_{0,4}^{\prime0}$ group of connected components of
\[
A'_{0,4}\subset\operatorname{Auttop}
\left(\deleted{\mathbb P^1_{\mathbb C}\setminus\unsure{\{}}
\underset{\substack{\Vert\\\mathbb P^1_{\mathbb C}\setminus\{0,1,\infty;\tfrac12\}}}{U_{0,4}^{\mathrm{top}}}\right)
\]
subgroup of automorphisms that fix the point at infinity $a_{0,4}=\tfrac12$

with $A_{0,4}^{\prime+}/A_{0,4}^{\prime0}$ (preserving orientation

4) $\Pi_1\bigl((U_{0,3}^{\mathrm{an}},\overbrace{\mathfrak S_3\times\unsure{\mathbb Z_2|\mathbb Z/2}}^{\mathcal{T}_{0,3}}),-\bar j\bigr)$
(described in the manner of J. Malgoire)
with the subgroup $\Pi_1\bigl((U_{0,3}^{\mathrm{an}},\mathfrak S_3),-\bar j\bigr)$
$(\mathfrak S_3=\mathcal{T}_{0,3}^{+})$

5) $\Pi_1\bigl((\mathbb P^{1\,\mathrm{an}}_{\mathbb C},0\cdot\{0\}+2\{1\}+3\{\infty\}),-\bar j\bigr)$
(in the manner of J. Malgoire)
with the subgroup $\Pi_1\bigl(\mathbb P^{1\,\mathrm{an}}_{\mathbb C},\unsure{0}\{0\}+2\{1\}+3\{\infty\}\bigr)$

6) Triangulated cartographic group \deleted{\illegible{1}} $\Gamma_{\infty,3}$
\[
\Gamma_{\infty,3}\unsure{\}}=\{r_0,r_1,r_\infty\mid
r_0^2=r_1^2=r_\infty^2=(r_0r_\infty)^2=(r_1r_0)^3=1\}
\]
(with oriented triangulated cartographic subgroup
\[
\Gamma_{\infty,3}^{\circ}=\{
\underset{\substack{\Vert\\r_\infty r_1}}{\rho_s},
\underset{\substack{\Vert\\r_0r_\infty}}{\sigma},
\underset{\substack{\Vert\\r_1r_0}}{\rho_f}
\mid\rho_f\sigma\rho_s=\sigma^2=\rho_f^3=1\}.
\]
)

\margin{this is also the group
$\widetilde{\Pi}_{0,3}/(\ell_1^2,\ell_\infty^3)$, where
$\widetilde{\Pi}_{0,3}=\Pi_{0,3}\cdot\{1,\tau\}$}

7) $\Pi_1\bigl(\mathcal U_{0,3}^{\mathrm{conf}},\operatorname{pt}(U_{0,3},-\bar j)\bigr)$
\inserted{universal conformal surface of type $(0,3)$}
with subgroup $\Pi_1\bigl(\mathcal U_{0,3}^{\mathrm{an}},\operatorname{pt}(U_{0,3},-\bar j)\bigr)$
\inserted{universal holomorphic surface of type $(0,3)$}

\margin{\uline{N.B.} $\mathcal U_{0,3}^{\mathrm{conf}}\simeq(U_{0,3},\mathcal T_{0,3})$
which makes it possible to identify 7) with 4)}

8) $\mathcal T'_{1,1}=\mathcal T_{1,1}/\underbrace{\{\pm1\}}_{\text{center}}$,
with $\mathcal T_{1,1}\simeq\operatorname{Autext}_{\mathrm{lac}}(\Pi_{1,1})
\simeq A_{1,1}/A_{1,1}^{0}$

($A_{1,1}\simeq\text{topological automorphisms of }\unsure{U_{1,1}^{\circ}|U_{1,1}}$)
with subgroup $\mathcal T_{1,1}^{\prime+}=\mathcal T_{1,1}^{+}/\{\pm1\}$

\scanpage{96}
\mspage{337}

9) $\Pi_1(M_{1,1}^{\mathrm{conf}},\mathrm{pt})' \simeq \unsure{\Pi_1()/\mathrm{centre}\{\pm1\}} \simeq \Pi_1(M_{1,1}'{}^{\mathrm{conf}},\mathrm{pt})$

where $M_{1,1}^{\mathrm{an}}$ is the \deleted{analytic}\inserted{\unsure{conformal|conformal}} modular multiplicity of \inserted{conformal} elliptic curves (i.e. \deleted{analytic}\inserted{conformal} curves of type $1,1$), $M_{1,1}'{}^{\mathrm{an}}$ that of \deleted{analytic}\inserted{conformal} elliptic curves ``modulo symmetry''; with subgroup
\[
\Pi_1(M_{1,1}^{\mathrm{an}},\mathrm{pt})' \simeq \unsure{\Pi_1()/\mathrm{centre}\{\pm1\}} \simeq \Pi_1(M_{1,1}'{}^{\mathrm{an}},\mathrm{pt})
\]
where the base point is the \inserted{conformal elliptic \unsure{resp.}} analytic curve, quotient of $\mathbb C$ by $\mathbb Z+\mathbb Z(-\bar j)$. \uline{N.B.} The modular multiplicities are $(D,\mathrm{Gl}(2,\mathbb Z))$, $(D,\mathrm{Sl}(2,\mathbb Z))$ in the \deleted{analytic} conformal resp. analytic setting for $M_{1,1}$, and $(D,\mathrm{Gl}(2,\mathbb Z)'=\mathrm{GP}(1,\mathbb Z))$ resp. $(D,\unsure{\mathrm{Sl}(1,\mathbb Z)'|\mathrm{Sl}(2,\mathbb Z)'}=\mathrm{GP}(1,\mathbb Z)^+)$ in the conformal resp. analytic setting for $M_{1,1}'$.

10) $\mathrm{Gl}(2,\unsure{\mathbb Z^{\ast}|\mathbb Z})' \simeq \mathrm{GP}(1,\mathbb Z)$, with the subgroup $\mathrm{Sl}(2,\mathbb Z)'\simeq\mathrm{GP}(1,\mathbb Z)^+$.

\scanpage{97}

(VI) \uline{The loop automorphisms of} $\widehat{\Pi}_{0,3}$ \uline{and of} $\widehat{\mathfrak S}_{0,3}^{+}$

We have a canonical \unsure{commutative} diagram:

(49)
\diagram{\begin{tikzcd}[ampersand replacement=\&,column sep=large,row sep=large]
\begin{array}{l}
\operatorname{Autext}(\widehat{\mathfrak S}_{0,3}^{+})'\\
=\text{group of outer automorphisms}\\
\text{inducing}\\
\text{the identity in the}\\
\text{quotient }\mathfrak S_3\simeq\mathcal T_{0,3}^{+}
\end{array}
\arrow[r,"{\text{restriction}}","\sim"']
\&
\begin{array}{l}
\operatorname{Autext}^{\unsure{\rho}}(\widehat{\Pi}_{0,3})'\\
=\text{group of outer automorphisms}\\
\text{that}\\
\text{commute with}\\
\text{the outer action of }\mathfrak S_3
\end{array}
\\
\begin{array}{l}
\operatorname{Aut}(\widehat{\mathfrak S}_{0,3}^{+},\sigma_\infty)'\\
=\text{group of automorphisms}\\
\text{of }\widehat{\mathfrak S}_{0,3}^{+}\text{ that }\deleted{\unsure{\text{fix}}}\inserted{\text{fix}}\ \sigma_\infty,\\
\text{and that fix }\rho\ \mathrm{mod}\ \widehat{\Pi}_{0,3}
\end{array}
\arrow[u,"\sim"]\arrow[r,"{\text{restriction}}","\sim"']
\&
\begin{array}{l}
\operatorname{Aut}(\widehat{\Pi}_{0,3},\sigma_\infty)\\
=\text{group of automorphisms}\\
\text{that commute}\\
\text{with }\sigma_\infty,\text{ and that}\\
\text{commute}\\
\mathrm{ext}^{\mathrm t}.\text{ with }\rho
\end{array}
\arrow[u,"\sim"]
\end{tikzcd}}

\margin{\deleted{\illegible{1}}}

The \unsure{$\in$} \deleted{\unsure{elements}} $\operatorname{Aut}(\widehat{\mathfrak S}_{0,3}^{+},\sigma_\infty)'$ are of the form
\[
(50)\qquad
\left\{\begin{aligned}
u(\sigma_\infty)&=\sigma_\infty\\
u(\rho)&=g\rho
\end{aligned}\right.
\qquad (\text{i.e. }u\rho u^{-1}\rho^{-1}=g)
\]

\margin{We write the formulas in $\operatorname{Aut}(\widehat{\mathfrak S}_{0,3}^{+})'=G$, group of automorphisms of $\widehat{\mathfrak S}_{0,3}^{+}$ that induce the identity in $\widehat{\mathfrak S}_{0,3}^{+}/\Pi_{0,3}\simeq\mathfrak S_3$, \uline{\unsure{is identified}} with the group $\operatorname{Aut}(\Pi_{0,3})'=\mathrm{gr.}$ of automorphisms of $\Pi_{0,3}$ that commute \inserted{\unsure{\uline{outer}}} with $\mathfrak S_3$; -- we shall have $\Pi_{0,3}\subset\widehat{\mathfrak S}_{0,3}^{+}\subset\widehat{\mathfrak S}_{0,3}\subset\unsure{\mathfrak g.|g.}\,G$. \uline{N.B.} We assume here that the center of $\widehat{\mathfrak S}_{0,3}^{+}$ is trivial.}

$g\in\widehat{\Pi}_{0,3}$ \deleted{\illegible{1}} is subject to the cocycle condition
\[
(51)\qquad \boxed{g\rho(g)\rho^2(g)=1}
\]
which expresses the relation
\[
(52)\qquad (g\rho)^3=1\ .
\]

We thus find a $1-1$ correspondence between the elements of $\operatorname{Aut}(\widehat{\mathfrak S}_{0,3}^{+},\sigma_\infty)'$ and the $g\in\widehat{\Pi}_{0,3}$ \unsure{satisfying} (51), and such moreover that

(53) $\sigma_\infty,g\rho$ with the relations $\sigma_\infty^2=(g\rho)^3=1$ \unsure{present} $\widehat{\mathfrak S}_{0,3}^{+}$

\scanpage{98}
\mspage{338}

(which I am going to call the ``preliminary condition'').

\deleted{For \unsure{$u$ thus defined by a}}

For a $u\in G=\operatorname{Aut}(\widehat{\mathfrak S}_{0,3}^{+})'\unsure{!}$ (not necessarily fixing $\sigma_\infty$), we have equivalent conditions
\[
(54)\quad
\left\{\begin{aligned}
\text{(i)}\quad &u(\overset{\inserted{\varepsilon_0}}{\ell'_0})\in\text{conjugate of }\deleted{\unsure{\ell}}\varepsilon_0^{\widehat{\mathbb Z}}=L'_0\ (\text{in }\widehat{\mathfrak S}_{0,3}^{+})\\
\text{(ii)}\quad &u(\ell_0)\in\text{conjugate of }\ell_0^{\widehat{\mathbb Z}}=L_0\ (\text{in }\widehat{\Pi}_{0,3})\\
\text{(iii)}\quad &u(L_i)\text{ conjugate of }L_i\quad\text{for }i=0,1,\infty\\
\text{(iv)}\quad &\exists\ p\in\widehat{\mathbb Z}^{\ast},\ g_0\in\widehat{\Pi}_{0,3}\quad
\boxed{u(\varepsilon_0)=\operatorname{int}(g_0)\varepsilon_0^p}\\
\text{(v)}\quad &\exists\ p\in\widehat{\mathbb Z}^{\ast},\ g_0\in\widehat{\Pi}_{0,3}\quad
\boxed{u(\ell_0)=\operatorname{int}(g_0)(\ell_0^p)}
\end{aligned}\right.
\]

We shall then also have
\[
(55)\qquad
\begin{aligned}
\deleted{\unsure{u}}\quad u(\ell_1)&=\operatorname{int}(g_1)(\ell_1^p)&
g_1&=g\rho(g_0)\\
u(\ell_\infty)&=\operatorname{int}(g_\infty)(\ell_\infty^p)&
g_\infty&=g\rho(g_1)=g\rho(g)\rho^2(g_0)\ .
\end{aligned}
\]

The element \deleted{$p\in\widehat{\mathbb Z}^{\ast}$ \unsure{is}}
\[
(55)\qquad p\in\widehat{\mathbb Z}^{\ast}\qquad p=\chi(u)
\]
is well determined by $u$, hence by $g$\inserted{(if $u$ given by (50))}, we shall call it the \uline{multiplier} of $u$. The element \deleted{\unsure{$g_0$}}
\[
(56)\qquad g_0\in\widehat{\Pi}_{0,3}
\]
is determined a priori modulo $L_0=\ell_0^{\widehat{\mathbb Z}}$ (change
\[
(57)\qquad g_0\longmapsto g_0\ell_0^\alpha
\]
).

\margin{Notations $\operatorname{Autext}_{\mathrm{lac}}(\widehat{\mathfrak S}_{0,3}^{+})$, $\operatorname{Aut}_{\mathrm{lac}}(\widehat{\mathfrak S}_{0,3}^{+})$, $\operatorname{Aut}_{\mathrm{lac}}(\widehat{\mathfrak S}_{0,3}^{+})'$, $\operatorname{Aut}_{\mathrm{lac}}(\widehat{\mathfrak S}_{0,3}^{+},\sigma_\infty)'$, $\operatorname{Aut}_{\mathrm{lac}}(\widehat{\Pi}_{0,3},\sigma_\infty)\ldots$}

If the \inserted{equivalent} conditions (\unsure{54|55}) are satisfied, we say that $u$ respects the loop structure of $\widehat{\mathfrak S}_{0,3}^{+}$ (or of $\Pi_{0,3}$) -- it being understood that $\widehat{\mathfrak S}_{0,3}^{+}$ is equipped with the ``loop-group'' $\varepsilon_0^{\widehat{\mathbb Z}}\simeq L'_0$. We \unsure{shall restrict|restrict} ourselves from now on to such automorphisms. \uline{N.B.} These conditions depend only on the outer automorphism defined by $u$, in $\widehat{\Pi}_{0,3}$ or in $\widehat{\mathfrak S}_{0,3}^{+}$. We

\scanpage{99}

are going to \unsure{take an interest} in the canonical liftings of such outer automorphisms given by (49), (50), to actual automorphisms.

The defining formula (50) for $u$ involves only $g$\inserted{$\in\widehat{\Pi}_{0,3}$}, in principle $g_0\ \mathrm{mod}\ L_0$ in $\widehat{\Pi}_{0,3}/L_0$ is deduced from $g$. But in practice $g_0$ plays as important a role as $g$ -- we shall see that to a large extent it determines $g$ (formula (62)).

If there were no need to introduce $g_0$, to express the conditions (54), we could easily \uline{parametrize} the $g$ \unsure{satisfying} (51) -- hence corresponding to \deleted{\unsure{splittings of}}\inserted{liftings of} $\dot\rho\in\mathfrak S_3$ to elements of order $3$ of $\widehat{\mathfrak S}_{0,3}^{+}$ -- by noting that there are exactly two conjugacy classes of such liftings, (corresponding to the fixed points $-\bar j,-j$ of $\dot\rho$ operating on $U_{0,3}^{\mathrm{an}}$), which are parametrized by
\[
(58)\quad
\left\{\begin{aligned}
g&=\lambda_0^{\unsure{\sharp|\ast}}\rho(\lambda_0^{\unsure{\sharp|\ast}})^{-1}
=[\lambda_0^{\unsure{\sharp|\ast}},\rho]
&&\begin{gathered}\text{(principal \deleted{first}\inserted{class}, \unsure{corresponding|corresponding}}\\
\text{to the ``north pole'' }-\bar j)\end{gathered}\\
g&=\lambda_0^{\unsure{\sharp|\ast}\prime}\ell_1^{-1}\rho(\lambda_0^{\unsure{\sharp|\ast}\prime})^{-1}
&&\begin{gathered}\text{(adjoint class, corresponding to the}\\
\text{``south pole'' }-j)\end{gathered}
\end{aligned}\right.
\]

\margin{\unsure{These|The} parametric \unsure{representations} are \unsure{unambiguous|unique}.}

where $\lambda_0,\lambda_0^{\unsure{\sharp|\ast}\prime}\in\widehat{\Pi}_{0,3}$. We shall see below that these cases correspond respectively to the cases
\[
(59)\qquad p\equiv1\ (3)\qquad\text{and}\qquad p\equiv-1\ (3)\ .
\]

\scanpage{100}
\mspage{339}

But here the situation gets considerably more involved through the need to take account of the \deleted{\unsure{the|the}} conditions (54) (for example in the form (iv)) -- hence \inserted{that} of introducing $g_0,p$ jointly. It is then appropriate to introduce the quantities
\[
(60)\qquad \gamma=\frac{p-1}{2}\in\widehat{\mathbb Z}
\]
(N.B. since $p\in\widehat{\mathbb Z}^{\ast}$, its \unsure{class} in $\mathbb Z/2\mathbb Z$\inserted{$=\widehat{\mathbb Z}/2\widehat{\mathbb Z}$} is $1$, hence $p-1\in2\widehat{\mathbb Z}$, whence the possibility of finding $\gamma$),
\[
(61)\qquad h_0=g_0^{-1}\sigma_\infty(g_0)
\]

\margin{$h_0$, relative to the automorphism $\sigma_\infty$ of order $2$, \unsure{plays} a role symmetric \unsure{to} $\xi$ relative to the automorphism $\rho$ of order $3$ -- this does not show up \unsure{either} in the notation.}

\deleted{(we change our notation, previously we had taken $h_0=\sigma_\infty(g_0)^{-1}g_0$, so we take the inverse \unsure{with respect|p.v.} to the previous notation).} We check \unsure{indeed|thus} that $g$ is expressed by means of $g_0$ and of $\gamma$ (hence of $p$) by means of the formula (\uline{equivalent} to (54) (iv) or (v))
\[
(62)\quad
\boxed{g=\sigma_\infty(g_0)\ell_1^\gamma\rho(g_0)^{-1}}
\left(=(g_0\ell_0^{-\gamma})\left(\underbrace{\ell_0^\gamma h_0}_{\xi}\right)\rho(g_0\ell_0^{-\gamma})^{-1}\right)
\]

\margin{\uline{N.B.} This expression for $g$ is not \unsure{changed} if $g_0\longmapsto g_0\ell_0^\alpha$.}

Thus the $u$ that interest us are described by the pairs $(g_0,\gamma)$ $(g_0\in\widehat{\Pi}_{0,3}/L_0,\ \gamma\in\widehat{\mathbb Z})$ \unsure{satisfying} the conditions
\[
(63)\qquad 2\gamma+1\in\widehat{\mathbb Z}^{\ast}\qquad (\mathrm{NB}\quad 2\gamma+1=p=\chi(u))
\]
and condition (51), where $g$ is expressed as a function of $g_0$ and of $\gamma$ by (62). This \deleted{formula}\inserted{relation} (51) is in turn equivalent to the relation
\[
(64)\qquad\boxed{\xi\rho(\xi)\rho^2(\xi)=1}
\]
where
\[
(65)\quad\xi=\ell_0^\gamma h_0
\qquad\unsure{\deleted{\text{with}}|\text{with}}\ (66)\quad
\begin{aligned}
h_0&\overset{\mathrm{dfn}}{=}
g_0^{-1}\sigma_\infty(g_0)\\
&{}=[g_0^{-1},\sigma_\infty]
\end{aligned}
\]

\margin{\uline{N.B.} We \unsure{call} $h_0$ here what we previously denoted by \unsure{$h_0^{-1}$}.}

\scanpage{101}

Thus, the $u$ are \unsure{parametrized} by the pairs
$(g_0\in\widehat{\Pi}_{0,3}/L_0,\ \gamma\in\widehat{\mathbb Z})$, satisfying conditions
(63) and (64), which (by means of (65)) is written in explicit form
\[
\tag{\unsure{65bis}}
\boxed{
\bigl(\ell_0^\gamma\underbrace{g_0^{-1}\sigma_\infty(g_0)}_{h_0}\bigr)
\rho\bigl(\ell_0^\gamma\underbrace{g_0^{-1}\sigma_\infty(g_0)}_{h_0}\bigr)
\rho^2\bigl(\ell_0^\gamma\underbrace{g_0^{-1}\sigma_\infty(g_0)}_{h_0}\bigr)=1
}
\]
\margin{it would be better to regard this equation as an ``equation in $h_0$'' \unsure{playing} a role symmetric to ``the equation in $\xi$'' (68)}
(``equation in $g_0$'' --- as opposed to the ``equation in $\xi$'' below).
(Moreover, we \unsure{must} add the preliminary condition, for the record...)

We can also consider that \deleted{$g$ is} $u$ (via $g$) is defined by
$(\xi\in\widehat{\Pi}_{0,3},\ \gamma\in\widehat{\mathbb Z})$, which indeed allow us to recover
\[
h_0=\ell_0^{-\gamma}\xi
\]
by (65), then $g_0$ by (66)
\[
g_0^{-1}\sigma_\infty(g_0)=h_0
\]
which indeed determines $g_0$ uniquely, thanks to the fact (which we admit provisionally)
that $\widehat{\Pi}_{0,3}^{\sigma_\infty}=\{1\}$.
The possibility of finding $g_0$ satisfying this condition is equivalent (we shall admit this here too) to the condition
\[
\tag{67}
\sigma_\infty(h_0)h_0=1
\]
which, in terms of $\xi,\gamma$, is also written in the \unsure{form}

\scanpage{102}
\mspage{340}

\[
\tag{68}
\boxed{\sigma_\infty(\xi)=\ell_1^\gamma\xi^{-1}\ell_0^\gamma}.
\]

Equations (64), (68) are called the ``equations in $\xi$'' --- as opposed to the ``equation in $g_0$'' (66), which at first sight has the advantage of being a single one, and the disadvantage of being more \unsure{complicated}.
\margin{rather ``the equation in $h_0$''}
In summary, the $u$ that interest us
$(u\in\mathrm{Aut}_{\mathrm{lac}}(\widehat{\mathfrak S}_{0,3}^{+},\sigma_\infty)')$
are \unsure{given|described} by the pairs
\deleted{$\unsure{(\xi,\gamma)}$}\inserted{$(\xi\in\widehat{\Pi}_{0,3}^{+},\ \gamma\in\widehat{\mathbb Z})$}, satisfying the conditions
\[
\begin{aligned}
\text{a)}\quad&2\gamma+1\in\widehat{\mathbb Z}^{*} &&(63)\\
\text{b)}\quad&\xi\rho(\xi)\rho^2(\xi)=1 &&(64)\\
\text{c)}\quad&\sigma_\infty(\xi)=\ell_1^\gamma\xi^{-1}\ell_0^\gamma &&(68)
\end{aligned}
\]
d) Preliminary condition for the record (53)
(which apparently are not \unsure{expressed} directly in a convenient way in terms of $\xi,\gamma$).

For two pairs $(\xi,\unsure{\gamma})$, $(\xi',\unsure{\gamma}')$ to define the same $u$, it is necessary and sufficient that we have
\[
\tag{69}
\unsure{\gamma}=\unsure{\gamma}',\qquad
\xi'=\ell_0^{-\alpha}\xi \ell_1^\alpha\qquad(\alpha\in\widehat{\mathbb Z}).
\]

Condition (64) can also be \unsure{replaced} by a parametric representation
\[
\tag{70}
\left\{
\begin{aligned}
\xi&=\lambda\rho(\lambda)^{-1}\qquad\text{or}\\
\xi&=\lambda \ell_0^{-1}\rho(\lambda)^{-1}
\end{aligned}
\right.
\]
with $\lambda\in\widehat{\Pi}_{0,3}$ --- in which case equation (68)

\scanpage{103}

can be interpreted as an ``equation in $\lambda$'', which according to the two cases takes the form
\[
\tag{71}
\left\{
\begin{aligned}
\sigma_\infty(\lambda\rho(\lambda)^{-1})
 &=\ell_1^\gamma(\rho(\lambda)\lambda^{-1})\ell_0^\gamma\\
&\text{or}\\
\sigma_\infty(\lambda\unsure{\ell_0^{-1}}\rho(\lambda)^{-1})
 &=\ell_1^\gamma(\rho(\lambda)\ell_0\lambda^{-1})\ell_0^\gamma,
\end{aligned}
\right.
\]

Here for $(\lambda,\gamma)$ and $(\lambda',\gamma')$ to define the same $u$, it is necessary and sufficient that we have \deleted{\illegible{1}}
\[
\tag{72}
\lambda'=\ell_0^{-\alpha}\lambda,\qquad\gamma=\gamma'.
\]

\uline{Remarks} 1) Up to the ``harmless'' factor $\ell_0^\gamma$, \unsure{which|that} distinguishes $h_0$ and $\xi$ (65), we can say that the \unsure{search} for the
$u\in\mathrm{Aut}_{\mathrm{lac}}(\widehat{\mathfrak S}_{0,3}^{+},\sigma_\infty)$
is equivalent to that for the \deleted{\unsure{equations}} \uline{simultaneous} solutions of the two cocycle equations
\[
\tag{73}
\left\{
\begin{aligned}
\xi\rho(\xi)\rho^2(\xi)&=1\\
h_0\sigma_\infty(h_0)&=1
\end{aligned}
\right.
\]
relative \deleted{respectively} to the groups $\mathbb Z/3\mathbb Z$, operating on $\widehat{\Pi}_{0,3}$ via $\rho$, and $\mathbb Z/2\mathbb Z$, operating via $\sigma_\infty$ --- whose solutions in parametric forms are respectively
\[
\tag{74}
\left\{
\begin{aligned}
\xi&=\lambda\rho(\lambda)^{-1}\quad\text{or}\quad
 \xi=\lambda \ell_0^{-1}\rho(\lambda)^{-1}\\
&\quad\text{(according as }\gamma\equiv0\ (3)\text{ or }\gamma\equiv\unsure{1|-1}\ (3)\\
&\qquad\text{i.e. }p\equiv1\ (3)\text{ or }p\equiv-1\ (3)\text{)}\\
h_0&\unsure{=|\equiv}\unsure{g_0^{-1}\sigma_\infty(g_0)}
\end{aligned}
\right.
\]
\margin{N.B. expression (62) for $g$ brings out equally well \unsure{this double} \unsure{intervention} of $\sigma_\infty$ and of $\rho$ in the determination of the $u$!!!}
\margin{(N.B. we could put \unsure{everywhere} $g_0$ for $g_0^{-1}$...)}

\scanpage{104}
\mspage{341}

3) The relation between the $\lambda_0$ in the parametric representation of $g$ (58), and the $\lambda$ in that of $\xi$ (74 a) is implicit in the last expression (62) for $g$, namely
\[
\tag{75}
\lambda_0=g_0\ell_0^{-\gamma}\lambda.
\]

4) Let us note here also, that the quantities
\[
\begin{array}{cccccccc}
p,&\gamma,&g,&g_0,&h_0,&\xi,&\lambda_0,&\lambda\\
\rotatebox{-90}{$\in$}&\rotatebox{-90}{$\in$}&\rotatebox{-90}{$\in$}&\rotatebox{-90}{$\in$}&\rotatebox{-90}{$\in$}&\rotatebox{-90}{$\in$}&\rotatebox{-90}{$\in$}&\rotatebox{-90}{$\in$}\\
\widehat{\mathbb Z}^{*}&\widehat{\mathbb Z}&\widehat{\Pi}_{0,3}&\widehat{\Pi}_{0,3}/L_0&&&\widehat{\Pi}_{0,3}&
\end{array}
\]
\inserted{resp. $\lambda'_0,\lambda'$}
(with \inserted{for $u$ constant} the indeterminacies specific to $g_0,h_0,\xi,\lambda$) can be considered \inserted{for $u$ variable,} as functions on the group
$\mathrm{Autext}_{\mathrm{lac}}(\widehat{\mathfrak S}_{0,3}^{+})'\simeq\mathrm{Autext}_{\mathrm{lac}}(\widehat{\Pi}_{0,3})'$
(subgroup of the outer automorphisms \unsure{of the group} with loops $\widehat{\Pi}_{0,3}$ that commute with $\mathfrak S_3$), with values in the sets
$\widehat{\mathbb Z}^{*},\widehat{\mathbb Z},\widehat{\Pi}_{0,3},\widehat{\Pi}_{0,3}/L_0,\ ,\ ,\widehat{\Pi}_{0,3},\ $.
In particular, we thus find 8 functions \inserted{\unsure{continuous}} on $\Gamma_{\mathbb Q}$, with values in the appropriate spaces --- starting with the cyclotomic character.
\margin{We shall see (\illegible{2} \unsure{less}!) that we can \uline{normalize} $g_0$, and thereby also $h_0,\xi,\lambda,\lambda_0$, so as to remove all ambiguity in \unsure{its|these} \unsure{definition|definitions}...}

\scanpage{105}

Let us note in this connection the equivalent conditions
\[
\tag{75bis}
\left\{
\begin{aligned}
\text{(i)}\quad &p=1\\
\text{(ii)}\quad &\gamma=0\\
\text{(iii)}\quad &g=\sigma_\infty(g_0)\rho(g_0)^{-1}\\
\text{(iv)}\quad &h_0=\xi\\
\text{(v)}\quad &\sigma_\infty(\xi)=\xi^{-1}
\qquad\text{(``equation in $\xi$'' simplified!)}\\
\text{(vi)}\quad &h_0\rho(h_0)\rho^2(h_0)=1
\end{aligned}
\right.
\]
all meaning that $\chi(u)=1$. Thus the search for the
$u\in\mathrm{Aut}_{\mathrm{lac}}(\widehat{\mathfrak S}_{0,3}^{+},\sigma_\infty)'$
\uline{with multiplier 1}
($u\in\mathrm{Aut}_{\mathrm{lac}}(\widehat{\mathfrak S}_{0,3}^{+},\sigma_\infty)^{\prime0}$)
is thus indeed equivalent to that for the simultaneous solutions of the equations
\inserted{of cocycles} (73), or of the parametric equations (``coboundary'') (74).

2) Returning to the general case (arbitrary multiplier), the only case that I know how to write down explicitly
\inserted{(except for $u=1$)} \unsure{is} the one that corresponds to
$\tau_\infty\in\unsure{\mathfrak S_{0,3}^{*}}$,
which corresponds to the characteristics
\[
\tag{75ter}
\left\{
\begin{aligned}
&u=\tau_\infty,\quad\chi(u)=p=-1,\quad\gamma=-1\\
&g=\ell_1^{-1}\\
&g_0=1,\quad h_0=1,\quad\xi=\ell_0^{-1}\\
&\lambda'_0=1,\quad\lambda'=\unsure{\ell_0^\gamma g_0^{-1}\lambda_0}=\ell_0^{-1}\\
&\inserted{[\,\lambda_0=\tau_\infty,\quad\lambda=\ell_0^{-1}\tau_\infty\,]}
\end{aligned}
\right.
\]
\margin{$\leftarrow$ cf. below}

It would certainly be no luxury \unsure{to manage} to write down a bunch of other $u$, to see in particular whether the \unsure{crowds|crowd} of necessary conditions that we are led to derive for the elements coming from $\unsure{\Gamma_{\mathbb Q}}$, are not automatically satisfied!

\scanpage{106}
\mspage{342}

(VII) Calculations in $\mathrm{Sl}(2,\widehat{\mathbb Z})'=\mathrm{Sl}(2,\widehat{\mathbb Z})/\pm1$
(\unsure{cocycles} under $\rho$).

\margin{Calculations ``for conscience's sake'', \unsure{which will not be of use} perhaps any more!}

This group is a (very small!) quotient of $\widehat{\mathrm{Sl}(2,\mathbb Z)'}$, which has the advantage of \unsure{lending} itself to convenient explicit \inserted{algebraic} calculations, in terms of matrices. Let us denote, by abuse of language, by the same letters $g,g_0,h_0,\xi,\lambda_0,\lambda$ elements of \deleted{\illegible{1}} $\widehat{\Pi}_{0,3}$, \inserted{(or of $\widehat{\mathfrak S}_{0,3}^{+}$ more generally)} and their images in \inserted{$\Pi(2)'\subset$} $\mathrm{Sl}(2,\widehat{\mathbb Z})'$ \deleted{\illegible{1}} (or in $\mathrm{Gl}(2,\widehat{\mathbb Z})'\simeq\mathrm{Gl}(2,\widehat{\mathbb Z})/\{\pm1\}$). \deleted{We} \deleted{\unsure{shall set}} in general, let us consider a matrix
\[
(76)\qquad g=\begin{pmatrix}a&b\\c&d\end{pmatrix}\in\mathrm{Gl}(2,\widehat{\mathbb Z})
\]
we then have
\[
(77)\quad
\left\{\begin{aligned}
\deleted{\unsure{\rho}}\rho(g)&=\rho g\rho^{-1}=\begin{pmatrix}c+d&-c\\c+d-a-b&a-c\end{pmatrix},
&\inserted{\det(g)}\rho(g^{-1})&=\begin{pmatrix}a-c&c\\a+b-c-d&c+d\end{pmatrix}\\
\rho^2(g)&=\rho^{-1}g\rho=\deleted{\unsure{(}}\begin{pmatrix}d-b&a+b-c-d\\-b&a+b\end{pmatrix},
&\inserted{\det(g)}\rho^2(g^{-1})&=\begin{pmatrix}a+b&c+d-a-b\\\deleted{\unsure{c}}b&d-b\end{pmatrix}
\end{aligned}\right.
\]
\[
(78)\quad
\left\{\begin{aligned}
\relax[g,\rho]&=g\rho g^{-1}\rho^{-1}=\begin{pmatrix}A&B\\C&D\end{pmatrix}.\ \deleted{\unsure{\frac{1}{\det(g)}}}\\
\delta A&=(a^2+b^2+ab)-(ac+bd+bc)&\delta B&=ac+bd+bc\\
\delta C&=-(c^2+d^2+cd)+(ac+bd+ad)&\delta D&=c^2+d^2+cd
\end{aligned}\right.
\]
\margin{N.B. $C+D-B=\unsure{1}\deleted{\unsure{\det g=\delta}}$\quad $\delta=\det g=ad-bc$ \deleted{\illegible{1}}}

We shall write the equation
\[
\deleted{(\unsure{79})}\qquad g\rho(g)\rho^2(g)=1
\]
in the equivalent form
\[
(79)\qquad g\rho(g)=\rho^{-1}(g^{-1})\bigl(=\rho^2(g^{-1})\bigr)
\qquad\text{\uline{in $\mathrm{Gl}(2,\widehat{\mathbb Z})'$}}
\]
where
\[
(80)\qquad g\rho(g)=\begin{pmatrix}P&Q\\R&S\end{pmatrix}
\]

\scanpage{107}

\[
(81)\quad
\left\{\begin{aligned}
P&=(a+b)(c+d-b)&Q&=-a(c-b)-bc\overset{\inserted{\deleted{\text{if }\det g=+1}}}{=}\deleted{\unsure{1}}\delta-a(c+d-b)\\
R&=(c+d)^2-d(a+b)&S&=ad-c(c+d)=\underset{\inserted{\text{if }\det g=1}}{\delta}-c(c+d-b)\\
&\underset{\inserted{\text{if }\det g=+1}}{=}(c+d)(c+d-b)-\delta
\end{aligned}\right.
\]
where we have set $\delta=ad-bc=\det g$.

Comparing (81) and the expression (77) for $\rho^2(g^{-1})$, equation (79) -- which is also written
\[
g\rho(g)=\pm\rho^2(g^{-1})
\]
gives the following solutions.

\margin{Up to replacing $g$ by $\unsure{-g}$, \deleted{\illegible{5}} \unsure{we are reduced} to \unsure{determining} the solutions of the equation}

\[
(82)\qquad \unsure{g|\rho}\rho(g)\rho^2(g)=\unsure{+|\pm}1\qquad\text{i.e.}\quad \unsure{g|\rho}\rho(g)=\unsure{+|\pm}\rho^{-1}(g^{-1})\qquad\text{\uline{in $\mathrm{Gl}(2,\widehat{\mathbb Z})$}}
\]
\deleted{$\unsure{(82)}$ \illegible{2} $\unsure{g\rho(g)=+\rho^{-1}(g^{-1})}$.}
We distinguish the solutions according to the factors of the product
\[
\mathrm{Sl}(2,\widehat{\mathbb Z})'\simeq\prod_{\ell\ \mathrm{premier}}\mathrm{Sl}(2,\mathbb Z_\ell)'.
\]
\margin{\uline{N.B.} \unsure{We} \unsure{shall use|use} this study of equation (82) only in the case $\det g=1$, i.e. $\delta=1$, \unsure{unless I am mistaken}}

The relation
\[
\frac1\delta(a+b)=P=(a+b)(c+d-b)\qquad\text{in }\mathbb Z_\ell
\]
gives the two solutions $a+b=0$ or $c+d-b=1$ (since $\mathbb Z_\ell$ is an integral domain).

\deleted{1°) $(a+b)\ne0$ (in $\mathbb Z_\ell$) whence $c+d-b=1$
\[
(83)\qquad\boxed{c+d-b=\frac1\delta}.
\]
Given (83), we verify \unsure{at once} that the formulas \inserted{$g\rho(g)=\rho^{-1}(g^{-1})$} -- which remain to be checked
\[
\begin{aligned}
P&=(a+b)/\delta&Q&=(1-a)/\delta\\
R&=b/\delta&S&=(1-c)/\delta
\end{aligned}
\]
\inserted{\deleted{\illegible{1}} $\unsure{R=}$
\[
\begin{aligned}
Q&=\frac{c+d-\unsure{a}-b}{\delta}\left(=\unsure{\frac1{\delta^2}|\frac1\delta}-\frac a\delta\right),\\
R&=\frac b\delta,\quad S=\frac{d-b}{\delta}\left(=\frac1{\delta^2}-\frac c\delta\right)
\end{aligned}
\]}
(84) are equivalent to \illegible{5} (82) \illegible{1}
\[
(84)\qquad\boxed{\delta^3=1}.
\]
}

\margin{The necessity of this formula (84) was obvious a priori.}

\deleted{2°) $a+b=0$ i.e. $a=-b$, i.e. $g=\begin{pmatrix}-b&b\\c&d\end{pmatrix}$, and
\[
-b(c+d)=1
\]
\[
(84)\quad\left\{\begin{aligned}
\det g&=-b(c+d)=\delta\quad(\text{thus }\delta+b(c+d)=0;\unsure{)}\quad d=-\frac\delta b-c\\
P&=0,&Q&=\delta+b(c+d-b)=-b^2\\
R&=(c+d)^2=\frac{\delta^2}{b^2},&S&=\unsure{ad}\ \illegible{1}\ {-c(c+d)+ad=1-c(\unsure{c+d-b})}\\
&&&\deleted{\illegible{1}c}=\delta+c\left(\frac\delta b+b\right)
\end{aligned}\right.
\]
}

\scanpage{108}
\mspage{343}

1°) $a+b\ne0$ (in $\mathbb Z_\ell$) whence
\[
(83)\qquad\boxed{c+d-b=\frac1\delta}.
\]
Given (83), we verify that the formulas (82), -- \unsure{where} it remains to check
\[
\begin{aligned}
Q&=\frac{c+d-(a+b)}\delta\ \unsure{\leftarrow|=}\left(=\frac1{\delta^2}-\frac a\delta\right),&R&=\frac b\delta\\
S&=\frac{d-b}\delta.\ \left(=\frac1{\delta^2}-\frac c\delta\right)
\end{aligned}
\]
are equivalent to:
\[
(84)\qquad\boxed{\delta^3=1}
\]
(whose necessity was obvious a priori).

\margin{\uline{N.B.} The conditions $c+d-b=\frac1\delta$, $\delta^3=1$ are sufficient \illegible{1} for equations (82), \unsure{!} (\unsure{independently} of the question whether $a+b=0$ or $a+b\ne0$) -- they are necessary \unsure{if we assume that we} \unsure{are not} in case 2° below}

2°) $a+b=0$ (in $\mathbb Z_\ell$) whence
\[
(85)\qquad a=-b\qquad\text{i.e.}\quad g=\begin{pmatrix}-b&b\\c&d\end{pmatrix}
\]
\[
(86)\quad\left\{\begin{aligned}
\det g&=-b(c+d)=\delta&\text{whence}\\
\deleted{\unsure{b+}\delta}\ \delta+b(c+d)&=0&d&=-\frac\delta b-c
\end{aligned}\right.
\]
i.e.
\[
g=\begin{pmatrix}-b&b\\c&-\frac\delta b-c\end{pmatrix},
\]
and
\[
\left\{\begin{aligned}
P&=0,&Q&=\delta+b(c+d-b)=-b^2\\
R&=(c+d)^2=\frac{\delta^2}{b^2},&S&=\delta+c\left(\frac\delta b+b\right);
\end{aligned}\right.
\]
\deleted{Equation (82) is then written, since \unsure{here}}
\[
\rho^2(g^{-1})=\begin{pmatrix}0&-1/b\\b/\delta&-\frac1\delta\left(\frac\delta b+b+c\right)\end{pmatrix},
\]
\deleted{and and}

and equation (82) is then written
\[
(87)\quad\left\{\begin{aligned}
\text{a)}\quad&b^3=\delta^3=1.\\
\text{b)}\quad&c\bigl((b\delta)^2+(b\delta)+1\bigr)+(\delta^2+b^2\delta+b)=0.
\end{aligned}\right.
\]

\scanpage{109}

We again distinguish two cases:

a) $b\delta=1$ i.e. $b=\frac1\delta$, then equation \inserted{87 b)} becomes
\[
3c+3\frac1\delta=0\qquad\text{i.e.}
\]
\deleted{$3c=$}
\[
(88)\qquad c=-\frac1\delta=\inserted{-\delta^2}\ \deleted{\unsure{3c+(2\delta^2+1)=0}}\qquad\deleted{\unsure{3c+3\delta^2=0}\quad\text{i.e.}}
\]
\deleted{Now $\delta\in\widehat{\mathbb Z}^{*}\Rightarrow\delta^2\equiv1\ (3)$ thus $2\delta^2+1\equiv0\ (3)$, whence the existence of a solution in $\widehat{\mathbb Z}$ of (88)
\[
(89)\qquad c=-\frac13(2\delta^2+1)
\]
}
corresponds to
\[
(90)\qquad g=\deleted{\begin{pmatrix}
-\frac1\delta=-\delta^2&\frac1\delta=\delta^2\\
\unsure{-\frac13(2\delta^2+1)}\ \inserted{-\delta^2}&\frac13(\unsure{4|1}-\delta^2)
\end{pmatrix}}\quad
\begin{pmatrix}
-\delta'\deleted{\unsure{{}^{2}}}&\delta'\deleted{\unsure{{}^{2}}}\\
-\delta'\deleted{\unsure{{}^{2}}}&0
\end{pmatrix}
\]
\inserted{where $\delta'=1/\delta=\delta^2$, and} where, I recall, $\delta$ is subject to the sole condition
\[
\delta^3=1.
\]

b) $b\delta\ne1$ thus (since $(b\delta)^3=1$) $(b\delta)^2+b\delta+1=0$)
and equation (87 b) becomes
\[
(\alpha)\qquad\delta^2+b^2\delta+b=0.\qquad\deleted{\text{\unsure{i.e. a}}}
\]
We shall have
\[
\left\{\begin{aligned}
b&=\deleted{\unsure{\delta}}u\frac1\delta=u\delta^2\quad\text{with}\quad u^3=1,\ u\ne1\quad\deleted{\text{thus}}\quad\text{i.e.}\\
u^2+u+1&=0
\end{aligned}\right.
\]
and equation $\alpha)$, \deleted{is written} reduces to
\[
\unsure{u^2|\delta^2}+u\deleted{\unsure{{}^{2}}}+1=0.
\]
\deleted{\illegible{1}}

\deleted{which requires $\delta\ne1$ (otherwise $u^2=0$, \deleted{\illegible{1}} incompatible with $u^3=1$) thus $\delta^2+\delta+1=0$ -- thus $(\alpha)$ is written
\[
-(\delta+1)-(u+1)+1=0\qquad\text{i.e.}
\]
\[
\delta+u+1=0
\]
which \deleted{is \unsure{obviously} \illegible{3}} \inserted{\unsure{means}} that $u,\delta$ are \uline{distinct} solutions of the cyclotomic equation $x^2+x+1=0$. Now $u\ne\delta$ \uline{means} \unsure{also that} $\unsure{\frac b\delta\ne\delta}$ i.e. $b\ne\delta^2$ \deleted{which requires} i.e. $b=\delta$.
\[
(91)\qquad g=\begin{pmatrix}-\delta&\delta\\c&-1-c\end{pmatrix}
\]
}

\scanpage{110}
\mspage{344}

\uline{Proposition} Let $g\in\mathrm{Gl}(2,\widehat{\mathbb Z}_\ell)$, $g=\begin{pmatrix}a&b\\c&d\end{pmatrix}$, $a,b,c,d\in\unsure{\widehat{\mathbb Z}_\ell|\mathbb Z_\ell}$,
$ad-bc=\delta\in\unsure{\widehat{\mathbb Z}_\ell^{*}|\mathbb Z_\ell^{*}}$. For us to have
\[
g\rho(g)\rho^2(g)=1\qquad(\unsure{82|81})
\]
it is necessary and sufficient that we have $\delta^3=1$, and that we be in one of the following two cases, which are mutually exclusive:
\[
\begin{aligned}
1\text{°)}\quad&c+d-b=1\\
2\text{°)}\quad&g=\deleted{\left(\begin{matrix}-\delta'&\delta'\end{matrix}\right.}\ -\delta'\begin{pmatrix}-1&1\\1&0\end{pmatrix}=-\delta'\rho^{-1}\qquad\left(\text{where }\delta'=\frac1\delta=\delta^2\right)
\end{aligned}
\]

\deleted{Moreover, in case 1°), we have $a+b=0$ \unsure{iff|sss} \unsure{when of}}

\uline{Corollary} The $g$ for which we have \deleted{at the same \uuline{time}} \inserted{1°) and} $a+b\ne0$ are those of the form
\[
\deleted{g=\left(\begin{matrix}-\unsure{u}\delta'&u\delta'\\\unsure{c|0}&\end{matrix}\right.}
\]
\[
\left\{\begin{aligned}
g&=\begin{pmatrix}a&b\\c&b-c+\frac1\delta\end{pmatrix}\qquad\deleted{\unsure{\text{for}}\ \unsure{c}=\frac1{(a+b)}\unsure{(a(b+\delta^2))|(ab+a\delta')}}\\
&\text{with}\quad c=\frac1{a+b}\bigl(a(b+\delta')-\delta\bigr),
\end{aligned}\right.
\]
those for which we have 1°) and $a+b=0$ are \deleted{those} the matrices of the form
\[
\left\{\begin{aligned}
g&=\begin{pmatrix}-u\delta'&u\delta'\\c&-\delta'u'-c\end{pmatrix}\\
&\text{where }\delta'=1/\delta=\delta^2,\ u'=\frac1u=u^2,\quad u\text{ being a solution}\\
&\text{of the cyclotomic equation }u^2+u+1=0\quad\text{i.e.}\\
&\text{\unsure{root|roots} \uline{\unsure{primitive|primitive}} }3^{\mathrm e}\text{ of unity}
\end{aligned}\right.
\]
